\documentclass[main.tex]{subfiles}
\begin{document}

\section{Additional Model Details}

\label{sec:oa-model}

\subsection{Consumer Payment Probabilities}\label{subsec:oa-payment-probs}

This section derives the payment probabilities $\pi^w_{M,j}$ from Section~\ref{subsubsec:consumption-over-merchants} for a model with two credit cards and one debit card, with results collected in Table~\ref{tab:consumer-payment-probs-baseline}.

Single-homers are straightforward: a consumer uses her card if the merchant accepts it and pays cash otherwise.
For multi-homers carrying two credit cards, substitution is also simple.
Both cards serve the same function, so if only one is accepted the consumer switches to it with probability~$1$ (middle panel of Table~\ref{tab:consumer-payment-probs-baseline}).

For a multi-homer carrying one credit and one debit card, credit and debit are segmented: each card is used only for its intended transaction type.
When both are accepted, the consumer uses credit with probability~$\pi$ and debit with probability~$1-\pi$.
When only one card is accepted, the consumer uses it for the share of transactions assigned to that type and pays cash for the rest.
A consumer whose credit card is accepted but whose debit card is not still uses credit with probability~$\pi$---she does not redirect debit transactions to the credit card.

\begin{table}[ht]
    \centering
    \caption{Consumer Payment Probabilities $\pi^w_M$ (Baseline)}
    \label{tab:consumer-payment-probs-baseline}
    \small
    \begin{tabular}{lll cc}
        \toprule
        \textbf{Primary Card} & \textbf{Secondary Card} & \textbf{Merchant Accepts} & $\pi^w_{M,w_1}$ & $\pi^w_{M,w_2}$ \\
        \midrule
        Any Card & None  & Primary card accepted      & 1 & 0 \\
        Any Card & None  & Primary card not accepted  & 0 & 0 \\
        \midrule
        Visa & MC  & Both Visa and MC  & $\pi$ & $1-\pi$ \\
        Visa & MC  & Only Visa         & 1 & 0 \\
        Visa & MC  & Only MC           & 0 & 1 \\
        Visa & MC  & Neither           & 0 & 0 \\
        \midrule
        Credit  & Debit & Both Credit and Debit  & $\pi$ & $1-\pi$ \\
        Credit  & Debit & Only Credit           & $\pi$ & 0 \\
        Credit  & Debit & Only Debit          & 0 & $1-\pi$ \\
        Credit  & Debit & Neither             & 0 & 0 \\
        \bottomrule
    \end{tabular}
\end{table}

\subsection{Acceptance Thresholds}\label{subsec:oa-credit-multihoming}

Every acceptance decision hinges on the same two forces: consumer multi-homing patterns and the fee spread across networks.
Consider two cards with fees $\tau_1 \leq \tau_2$.
Every payment method generates the same incremental sales per dollar.
A merchant already accepting card~$1$ decides whether to also accept card~$2$.
The incremental quasi-profit (Theorem~\ref{thm:quasiprofit-approx}) depends on two consumer groups:
\begin{itemize}
\item \textbf{New transactions} ($N$): single-homers on card~$2$ who would otherwise pay cash, generating incremental fees and sales.
\item \textbf{Diverted transactions} ($D$): multi-homers who switch from card~$1$ to card~$2$, changing the fee without generating new card transactions.
\end{itemize}
Both $N$ and $D$ are weighted sums of wallet shares $\mu^{(\cdot)}$, with weights depending on $\pi$ and card holdings.
The incremental quasi-profit coefficients are:
\begin{align}
\Delta a &= N\tau_2 + D(\tau_2 - \tau_1) \label{eq:a12-decomp}\\
\Delta b &= \tfrac{1}{\sigma}\bigl[N(1-\sigma\tau_2) - D\sigma(\tau_2-\tau_1)\bigr] \label{eq:b12-decomp}
\end{align}
The merchant accepts card~$2$ if and only if $\gamma \geq \underline{\gamma}_2$, where
\begin{equation}
\underline{\gamma}_2 = \frac{\sigma\bigl[N\tau_2 + D(\tau_2 - \tau_1)\bigr]}{N(1-\sigma\tau_2) - D\sigma(\tau_2-\tau_1)} \label{eq:general-accept-threshold}
\end{equation}
This threshold rule applies whenever $\Delta b > 0$, i.e., the incremental sales gain from accepting card~$2$ is positive; in the empirical setting $\Delta b > 0$ holds for all observed networks at baseline parameters. When $\Delta b \leq 0$, accepting the card destroys value for every type and the merchant declines unconditionally.

\subsubsection{Two Credit Cards}\label{subsubsec:oa-two-credit}

Multi-homing across credit cards affects acceptance, and relative fees matter.
Since both are credit cards, multi-homers switch to whichever is accepted (Table~\ref{tab:consumer-payment-probs-baseline}, middle panel), so all divert: $D = (1-\pi)\mu^{(1,2)} + \pi\mu^{(2,1)}$.
The new-transaction mass is $N = \mu^{(2,0)}$ (single-homers on card~$2$).
Substituting into~\eqref{eq:general-accept-threshold}:

\begin{equation}
\underline{\gamma}_2 = \frac{\sigma\mu^{(2,0)}\tau_2 + \sigma\left[(1-\pi)\mu^{(1,2)} + \pi\mu^{(2,1)}\right](\tau_2-\tau_1)}{\mu^{(2,0)}\left(1-\sigma\tau_2\right) - \sigma\left[(1-\pi)\mu^{(1,2)} + \pi\mu^{(2,1)}\right](\tau_2-\tau_1)} \label{eq:amex-accept-threshold}
\end{equation}

By inspection, cutting $\tau_1$ raises $\underline{\gamma}_2$ by widening $\tau_2 - \tau_1$, which increases diversion drag without changing incremental sales. Three special cases follow: $D = 0$ (no multi-homers) gives $\underline{\gamma}_2 = \sigma\tau_2/(1-\sigma\tau_2)$; $N \to 0$ (no single-homers) sends $\underline{\gamma}_2 \to \infty$; and $\tau_2 = \tau_1$ collapses the threshold to $\underline{\gamma}_1$.
This is consistent with Figures~\ref{fig:fees-history} and~\ref{fig:merchant-amex-visa}: AmEx acceptance responded sharply to its fee relative to Visa.


\subsubsection{Credit and Debit (Baseline)}\label{subsubsec:oa-credit-debit-baseline}

Multi-homing across credit and debit does not affect credit acceptance, and relative fees do not matter.
Credit and debit are segmented at the point of sale (Table~\ref{tab:consumer-payment-probs-baseline}, bottom panel): adding credit does not divert any spending from debit, so $D = 0$.
The threshold reduces to
\[
\underline{\gamma}_c = \frac{\sigma\tau_c}{1-\sigma\tau_c}
\]
independent of the debit fee.
Segmentation ensures $\Delta a$ and $\Delta b$ share the same wallet-share weights, so the threshold depends only on credit's own fee.
This is consistent with the Durbin evidence: after debit interchange fees fell by roughly half, credit interchange fees did not change (Figure~\ref{fig:merchant-fees-acceptance}) and yet credit card acceptance was unchanged (Figure~\ref{fig:scpc-debit-credit-acceptance}).

\subsection{CES Price Index and Multi-Homing Pecuniary Utility}\label{subsec:oa-ces-price-index}

A more natural pecuniary utility for multi-homers would use the CES price index directly:
\begin{equation*}
\log U^w = \pi f^{w_1} + (1-\pi) f^{w_2} - \log P^{(w_1,w_2)}.
\end{equation*}
I do not use the CES-index definition because it counterfactually predicts that Visa-AmEx multi-homing rates fall as AmEx acceptance rises (Appendix~\ref{subsec:oa-consumer-multihoming-history}).

When $w_1$ and $w_2$ are different types (credit and debit), the card-payment probability decomposes linearly: $v^w_M = \pi \cdot \mathbf{1}[w_1 \in M] + (1-\pi) \cdot \mathbf{1}[w_2 \in M]$, and $(P^w)^{1-\sigma} = \pi (P^{(w_1,0)})^{1-\sigma} + (1-\pi) (P^{(w_2,0)})^{1-\sigma}$.
The only remaining gap is a Jensen correction in the log price index, second-order in the difference between credit and debit price indices.

For same-type multi-homers (e.g., two credit cards), there is an additional first-order gap from union acceptance: the consumer pays by card whenever \emph{either} card is accepted, so $\pi \log P^{(w_1,0)} + (1 - \pi) \log P^{(w_2,0)} \geq \log P^{(w_1, w_2)}$, with equality when $w_1$ and $w_2$ have identical acceptance sets.
Since credit card acceptance is symmetric across networks in the estimated model and all counterfactuals, the same-type multi-homer first-order gap vanishes. 
For different-type multi-homers, the second-order Jensen wedge above remains.

\subsection{Details on the Conduct Assumption}\label{subsec:oa-conduct}

\paragraph{Why networks compete in gains $A_j$ rather than utilities or rewards.}
\textcite{Weyl2010} shows that guaranteeing utility gains captures penetration pricing when merchant acceptance is low, giving consumers a dominant adoption strategy and pinning down a unique subgame equilibrium.

\paragraph{Non-differentiability with multiple networks.}
Network profits are non-differentiable in merchant fees when more than two networks compete.
A network raising its fee competes with acceptance of all other networks, while a network cutting its fee competes with cash.
These regions have different marginal revenues, creating non-differentiability at the symmetric equilibrium.
\textcite{Rochet.Tirole2003}'s two-network model avoids this issue. \textcite{Teh.etal2022} show the problem arises with three or more networks.

\paragraph{Trembles for equilibrium selection.}
To select a unique equilibrium, I introduce small trembles $\nu_x$ in networks' fee choices, smoothing the profit function and ensuring equal profit loss from marginal fee increases or decreases (following \textcite{Bilal.Lhuillier2021}).
This convolution smooths the profit function, making $\Psi$ differentiable, and as the perturbation vanishes, the smoothed objective converges uniformly to the original.
To assess approximation error, Figure \ref{fig:max-compare} plots pairwise heatmaps of Visa's profit (without trembles) across all control-variable pairs, comparing the perturbed-FOC solution with a grid search over the original profit function.
The match is almost exact.

\begin{figure}[ht]
    \caption{Comparing the maximum from solving the FOCs of the perturbed profit function with the maximum of the original profit function}
    \label{fig:max-compare}
    \begin{center}
        \includegraphics*[width=5in]{../output/graphs/pairwise/baseline_1/baseline_visa}
    \end{center}
\end{figure}

\subsection{Model Solution Algorithm}
\label{subsec:oa-model-solution}

The algorithm has two layers.
An inner fixed point solves for the allocation (consumer wallet shares and merchant acceptance) given network fees $A_j$ and $\tau_j$.
An outer loop then jointly solves networks' first-order conditions over fees, taking the inner equilibrium as given.

\subsubsection{Fixed Point Iteration}

The inner fixed point iterates on wallet-specific price indices $P^w$.
Given a conjecture for these price indices, the algorithm recovers the implied uniform reward rates by inversion and proceeds through the following steps:
\begin{enumerate}
    \item Compute consumer wallet shares $\tilde{\mu}^w_y$ for each income level.
    \item Solve merchant acceptance: each merchant type $\gamma$ maximizes quasi-profit to determine the optimal card bundle $M^*(\gamma)$.
    \item Update price indices $P^w$ by integrating the price-index integrand over the merchant type distribution.
    \item Check convergence: if $\max_{w,s} |P^{w,\text{new}} - P^{w,\text{old}}| < \epsilon$, stop. Otherwise, apply damping and return to step 1.
\end{enumerate}
Profits for each network follow from Equations \ref{eq:transaction-fee-profit} and \ref{eq:subsidy-bill}.
Demand integration uses $2{,}000$ Monte Carlo draws jointly over unobserved preference heterogeneity and log-normal income.
% CONSTANT: simulation draw count, design choice

The outer loop solves networks' first-order conditions jointly, respecting the ownership matrix that accounts for Visa's co-ownership of credit and debit products.
Capped fees (e.g., under Durbin or counterfactual caps) bypass the first-order condition and are instead confirmed binding via Lagrange multiplier checks.
For the normal expectations in the trembles, $N$-dimensional Gauss-Hermite quadrature with 2 points per dimension ensures equal profit loss from raising or lowering fees.
% CONSTANT: 2 quadrature points per dimension, design choice
Second-order conditions for local maxima are verified numerically at every solution.

\subsection{Numerical Validation of the Quasi-Profit Approximation}
\label{subsec:oa-quasiprofit-validation}

The quasi-profit functions match true profits closely.
Figure \ref{fig:quasiprofit-approx} compares exact and approximate profits for all $32$ potential payment bundles in the baseline equilibrium.
% CONSTANT: 32 = 2^5 payment bundles from 5 payment methods
The fit holds across all values of $\gamma$.
Figure \ref{fig:quasiprofit-approx-envelope-example} plots the upper envelopes of true and approximate profit functions.
Slope changes in the envelope mark where the merchant's optimal acceptance strategy shifts.
The close fit confirms that the approximation accurately captures how optimal acceptance varies with merchant type.

\begin{figure}[ht!]
\caption{Comparison of exact profits (which factor in price changes) and quasi-profits (which do not) in the baseline equilibrium for every possible subset of accepted cards. \label{fig:quasiprofit-approx}}
\centering \includegraphics[width=0.8\linewidth]{../output/graphs/linear_approx_example.pdf}
\end{figure}

\begin{figure}[ht!]
    \caption{Comparison of the upper envelope of profits computed with the exact profit formula versus the quasi-profit approximation. \label{fig:quasiprofit-approx-envelope-example}}

    \centering \includegraphics[width=0.8\linewidth]{../output/graphs/linear_approx_envelope_example.pdf}
\end{figure}

\subsection{Mechanical Welfare Adjustment for Dual Routing}
\label{subsec:oa-mechanical-dr}

The dual routing counterfactual in Section~\ref{subsec:dual-routing} differs from the other counterfactuals because it directly perturbs preferences rather than fees, ownership, or the set of products. 
I raise $\chi^{22}$, the credit--credit complementarity parameter in the consumer's portfolio choice, by $\Delta \chi = 0.0020 \cdot \alpha_0$, equivalent to a 20 bp reward increase on credit--credit bundles for the median consumer.
% CONSTANT: Δχ is the policy magnitude defining this counterfactual scenario
A direct preference perturbation mechanically raises the consumer surplus in Equation~\ref{eq:cons-surplus-integral} through the logit log-sum, even if prices, portfolios, and merchant acceptance were held fixed.
To keep the dual routing results comparable with the other counterfactuals in the welfare tables, I subtract this mechanical component so that the reported consumer surplus reflects only the equilibrium response.

\paragraph{Income-specific adjustment.}
Let $\mathcal{W}_{CC} = \{w = (w_1, w_2) \in \mathcal{W} : w_1, w_2 \text{ are both credit cards}\}$ denote the set of credit--credit wallets.
For a consumer of type $i$ with reward-sensitivity $\alpha_i$, the boost $\Delta \chi$ on credit--credit wallets raises consumption-equivalent welfare, to first order, by
\begin{equation}
    w^{\text{mech}}_i(\tilde{\mu}) = \frac{\Delta \chi}{\alpha_i} \cdot P(\mathcal{W}_{CC}\mid i; \tilde{\mu}),
    \label{eq:mech-welfare-i}
\end{equation}
where $P(\mathcal{W}_{CC}\mid i; \tilde{\mu}) = \sum_{w\in\mathcal{W}_{CC}} \tilde{\mu}^{w}_i$ is $i$'s probability of holding a wallet in $\mathcal{W}_{CC}$, i.e., two credit cards.
By the Envelope theorem, evaluating $w^{\text{mech}}_i$ at the baseline portfolio shares $\tilde{\mu}^{\text{base}}$ is sufficient: the equilibrium change in shares enters welfare only at second order.
In practice, evaluating at counterfactual shares gives an answer indistinguishable from the baseline-shares value, consistent with this prediction.
Dividing by $\alpha_i$ converts utils to equivalent consumption, and weighting by $P(\mathcal{W}_{CC}\mid i; \tilde{\mu})$ accounts for the fact that only credit--credit wallets receive the boost.

The adjustment is larger for higher-income consumers: the consumption-equivalent adjustment averages 1.7 bps at low income, 3.9 bps at median, and 6.0 bps at high income in the baseline specification (Table~\ref{tab:mechanical-dual-routing-baseline}).
% HARDCODED[table: mechanical_dual_routing_baseline.tex, row: Low] current=1.7
% HARDCODED[table: mechanical_dual_routing_baseline.tex, row: Median] current=3.9
% HARDCODED[table: mechanical_dual_routing_baseline.tex, row: High] current=6.0
The pattern reflects higher credit--credit holding rates among high-income consumers, which more than offsets their higher reward-sensitivity.

\paragraph{Aggregate adjustment.}
The aggregate dollar gain scales the income-specific term \eqref{eq:mech-welfare-i} by income $y$ and integrates against the baseline income distribution $F$, mirroring Equation~\ref{eq:cons-surplus-integral}:
\begin{equation}
    W^{\text{mech}}(\tilde{\mu}) = \int \E\bigl[w^{\text{mech}}_i(\tilde{\mu}) \bigm| y\bigr] \cdot y \, dF(y).
    \label{eq:mech-welfare}
\end{equation}
Both integrals are computed by Monte Carlo over 500 joint draws of income and preference shocks per bootstrap sample.
% CONSTANT: 500 Monte Carlo draws for the mechanical welfare integral, design choice
The welfare tables subtract the mechanical effect on consumer surplus before reporting, so the \$3.3 billion consumer-surplus figure in Section~\ref{subsec:dual-routing} is net of this adjustment. The mechanical component is roughly \$4.5 billion in the baseline specification (Table~\ref{tab:mechanical-dual-routing-baseline}).
% HARDCODED[table: welfare_table_baseline.tex, row: "Consumers", col: Dual Routing] current=3.3
% HARDCODED[table: mechanical_dual_routing_baseline.tex, row: Aggregate] current=4.5

\begin{table}[ht!]
    \centering
    \caption{Mechanical dual routing welfare adjustment}
    \label{tab:mechanical-dual-routing-baseline}
    \input{../output/tables/mechanical_dual_routing_baseline.tex}
    \tablenote{Bootstrap standard errors in parentheses. The aggregate is in dollars, normalizing to \$10 trillion of total consumer purchases \parencite{Nilson2020a}. By-income values are the consumption-equivalent adjustment $w^{\text{mech}}_i$ in Equation~\eqref{eq:mech-welfare-i}, averaged across consumers within each income group, in basis points.}
\end{table}

\clearpage

\subsection{A Microfoundation for Segmentation Between Credit and Debit at the Point of Sale}
\label{subsec:oa-usage-microfoundation}

Payment methods can be substitutes at adoption yet poor substitutes at the point of sale.
A microfoundation building on \textcite{Koulayev.etal2016} and \textcite{Huynh.etal2022} separates usage from adoption to show that the two stages are independent.
It also gives a natural origin for the complementarity parameters $\chi^{lm}$ in the main text: carrying two different card types generates option value, while carrying two cards of the same type generates diminishing returns.
Rewards shift consumers between debit and credit at adoption even when the two are segmented at usage.

\subsubsection{Utility at the Point of Sale}

At the point of sale, the consumer picks one card from her wallet. Let $\mathcal{J} = \{0, 1, 2, \ldots, J\}$ denote the set of payment methods, where $j = 0$ is cash, and let $W$ denote the cards in consumer $i$'s wallet. Let $X_j \in \mathbb{R}^k$ be the characteristics of method $j$ (e.g., whether $j$ is a credit card), and $f_j$ the reward for using $j$.
For transaction $t$, the utility of method $j$ is:

\begin{align*}
u_{ijt} &= \alpha f_j + \gamma_{it}' X_j + \epsilon_{ijt} \\
\gamma_{it} &\sim \mathcal{N}(\gamma_i, \Sigma) \\
\epsilon_{ijt} &\sim \text{T1EV}
\end{align*}
The random shocks $\gamma_{it} \in \mathbb{R}^k$ capture transaction-specific preferences over card types: for a given transaction, a consumer may strongly prefer credit over debit, or cards over cash.
The mean $\gamma_i \in \mathbb{R}^k$ is consumer-specific, so some consumers prefer credit cards in general.
The idiosyncratic shocks $\epsilon_{ijt}$ capture why different cards of the same type (e.g., different credit cards) vary in value across transactions.

These assumptions yield the expected utility from wallet $W_i$:

\[
V_i(W_i) = \int \log\!\left(\sum_{j \in W_i} \exp\!\left(\alpha f_j + \gamma_{it}' X_j\right)\right) dH(\gamma_{it})
\]
where $H$ captures the distribution of $\gamma_{it}$.

\subsubsection{Utility at Adoption}

At adoption, consumers choose a wallet to maximize expected point-of-sale utility less adoption costs:

\begin{align*}
W_i^* &= \argmax_{W \subseteq \mathcal{J}} V_i(W) - \sum_{j \in W} c' X_j + \phi a_i^W \\
a_i^W &\sim \mathcal{N}(0, \nu^2)
\end{align*}
where $c \in \mathbb{R}^k$ captures adoption costs by payment-method characteristics.

\subsubsection{Mapping to the Model in the Main Text}

The microfounded model clarifies two points: how segmentation at the point of sale differs from segmentation at adoption, and how payment methods can complement or substitute for each other in a wallet.


\paragraph{Usage-stage vs.\ adoption-stage segmentation.} The main text assumes consumers use credit and debit for distinct transaction types, yet those carrying two credit cards freely substitute between them.
This is consistent with strong transaction-level preferences over credit versus debit but low heterogeneity in adoption costs across cards.
Strong usage-stage preferences arise when $\Sigma$ is large: a consumer wanting a credit feature for a given transaction gets little utility from a debit card.
This corresponds to certain transactions (e.g., large purchases) being better suited to credit than debit.

Strong adoption-stage substitution is compatible with $\nu^2 \to 0$.
Large $\Sigma$ does not induce segmentation at adoption because the consumer values the expected utility $V_i$, which integrates over the transaction-specific shocks $\gamma_{it}$.
If adoption costs are similar across instruments, consumers remain sensitive to rewards in choosing what to adopt.

As an extreme example, suppose a consumer has 50 large-ticket and 50 small-ticket transactions, strongly preferring credit for the former and debit for the latter.
At adoption, credit and debit offer similar expected utilities because each covers half of transactions.
Her adoption choice is therefore sensitive to rewards even though her transaction-level preferences are sharply bifurcated.

\paragraph{Complements and substitutes at adoption.} The microfounded model also explains why debit holders may want to adopt credit cards.
At adoption, consumers know their mean preferences $\gamma_i$ but have not yet drawn the transaction shocks $\gamma_{it}$; those shocks are realized only at the point of sale.
A consumer with high $\Sigma$ faces uncertainty over which card type she will prefer for each future transaction.
Adopting both types lets her use credit when $\gamma_{it}' X_j$ favors credit and debit when it favors debit, raising $V_i$.
This option value is positive even though the consumer knows her transaction-level preferences will be sharply bifurcated.

The model also captures incremental adoption costs.
If all cards of a given type pay the same rewards $f_j$, then $V_i$ is concave in the number of same-type cards, but costs are linear.
This explains why consumers do not adopt every card on the market, captured in the main text by $\chi^{22} < 0$ for wallets with two credit cards from different networks.

A consumer who highly values credit cards may still want a credit card alongside a debit card to capture more transaction-specific shocks $\gamma_{it}$, which occurs when $\gamma_i' X_j$ is large for credit.
The main text captures this through $\chi^{12} > 0$ and $\chi^{21} > 0$ for wallets that combine credit and debit.

\end{document}
